% =====================================================================
%  Module Data Lakehouse -- TP 3 : Images, orchestration et BI
%  Compilation : pdflatex (2 passes)
% =====================================================================
\documentclass[11pt,a4paper]{article}

\usepackage[utf8]{inputenc}
\usepackage[T1]{fontenc}
\usepackage[french]{babel}
\usepackage{lmodern}
\usepackage[babel=true]{microtype}

\usepackage[a4paper,margin=2.3cm,headheight=14pt]{geometry}
\usepackage{fancyhdr}
\usepackage{titlesec}
\usepackage{needspace}
\usepackage{enumitem}
\usepackage{amssymb}

\usepackage[table]{xcolor}
\usepackage{booktabs}
\usepackage{tabularx}
\usepackage{array}
\usepackage{tikz}
\usetikzlibrary{positioning,arrows.meta,babel}
\usepackage{listings}

\usepackage[htt]{hyphenat}
\usepackage[hidelinks]{hyperref}
\usepackage{xurl}
\usepackage{caption}
\usepackage{float}
\usepackage{graphicx}
% Captures d'écran : <nom>.png à la racine du projet (Overleaf) ou dans images/tp3/ ;
% tant que le fichier manque, un cadre indique le nom attendu.
\newcommand{\capture}[3]{%
  \begin{figure}[H]\centering
  \IfFileExists{#1.png}%
    {\includegraphics[width=0.92\textwidth,height=0.42\textheight,keepaspectratio]{#1.png}}%
    {\IfFileExists{images/tp3/#1.png}%
      {\includegraphics[width=0.92\textwidth,height=0.42\textheight,keepaspectratio]{images/tp3/#1.png}}%
      {\fbox{\parbox[c][3.2cm][c]{0.88\textwidth}{\centering\footnotesize
         Capture à ajouter : \texttt{\detokenize{#1}.png}}}}}%
  \caption{#2}\label{#3}
  \end{figure}}
\captionsetup{font=small,labelfont=bf,labelsep=period}

\setlength{\parskip}{0.5em}
\setlength{\parindent}{0pt}
\setlength{\emergencystretch}{3em}

% ---------- Couleurs ----------
\definecolor{lhblue}{HTML}{1F4E79}
\definecolor{lhteal}{HTML}{0F766E}
\definecolor{lhred}{HTML}{B91C1C}
\definecolor{lhviolet}{HTML}{6D28D9}
\definecolor{lhgray}{HTML}{F3F4F6}
\definecolor{lhbronze}{HTML}{A16207}
\definecolor{lhsilver}{HTML}{64748B}
\definecolor{lhgold}{HTML}{CA8A04}

% ---------- Titres, en-têtes ----------
\titleformat{\section}{\needspace{6\baselineskip}\Large\bfseries\color{lhblue}}{\thesection.}{0.6em}{}
\titleformat{\subsection}{\large\bfseries\color{lhblue}}{\thesubsection}{0.6em}{}
\pagestyle{fancy}
\fancyhf{}
\fancyhead[L]{\small\color{lhblue}Module Data Lakehouse}
\fancyhead[R]{\small\color{lhblue}TP 3 -- Images, orchestration et BI}
\fancyfoot[C]{\small\thepage}

% ---------- Encadrés : texte simple, sans boîte ----------
% Un titre en gras introduit le paragraphe ; seul « Attention » est en rouge.
\newenvironment{resultat}{\par\medskip\noindent\textbf{Résultat attendu :}\ }{\par\medskip}
\newenvironment{astuce}{\par\medskip\noindent\textbf{Astuce :}\ }{\par\medskip}
\newenvironment{attention}{\par\medskip\noindent\textbf{\textcolor{lhred}{Attention :}}\ }{\par\medskip}

% ---------- Code ----------
\lstset{
  basicstyle=\ttfamily\small, columns=fullflexible, keepspaces=true,
  breaklines=true, showstringspaces=false, upquote=true,
  frame=single, framesep=5pt, xleftmargin=6pt, xrightmargin=6pt,
  backgroundcolor=\color{lhgray}, rulecolor=\color{black!15},
  commentstyle=\color{black!55}\itshape,
  keywordstyle=\color{lhblue}\bfseries, stringstyle=\color{black!70},
  literate=
   {é}{{\'e}}1 {è}{{\`e}}1 {ê}{{\^e}}1 {à}{{\`a}}1 {ç}{{\c c}}1
   {ù}{{\`u}}1 {î}{{\^i}}1 {ô}{{\^o}}1 {É}{{\'E}}1 {Ê}{{\^E}}1
}
\lstdefinestyle{shell}{language=bash,backgroundcolor=\color{black!85},
  basicstyle=\ttfamily\small\color{white},rulecolor=\color{black!85},
  commentstyle=\color{black!35}\itshape,keywordstyle=\color{white}}

\newcommand{\cmd}[1]{\texttt{#1}}
\newcommand{\cocher}{$\square$\ }
\newcommand{\etape}[1]{\tikz[baseline=(n.base)]\node[circle,fill=lhblue,text=white,
  font=\sffamily\bfseries\scriptsize,inner sep=1.3pt](n){#1};}
\pdfstringdefDisableCommands{\def\cmd#1{#1}}

\newenvironment{notion}{\par\medskip\noindent\textbf{\textcolor{lhblue}{Notion clé :}}\ }{\par\medskip}
\newcommand{\nb}[1]{\textcolor{lhblue}{\texttt{#1}}}

% =====================================================================
\begin{document}
% =====================================================================

\begin{tikzpicture}[remember picture,overlay]
  \fill[lhblue] (current page.north west) rectangle ([yshift=-5.2cm]current page.north east);
  \node[anchor=north west,text=white,font=\sffamily] at ([xshift=2.3cm,yshift=-1.3cm]current page.north west)
    {Module Data Lakehouse -- Cycle ingénieur 4\textsuperscript{e}/5\textsuperscript{e} année};
  \node[anchor=north west,text=white,font=\sffamily\bfseries\fontsize{26}{30}\selectfont] at ([xshift=2.3cm,yshift=-2.2cm]current page.north west)
    {TP 3 -- Images, orchestration et BI};
  \node[anchor=north west,text=white!85!lhblue,font=\sffamily\large] at ([xshift=2.3cm,yshift=-3.6cm]current page.north west)
    {Un pipeline de bout en bout : automatisé, contrôlé, optimisé et exposé};
\end{tikzpicture}
\thispagestyle{empty}
\vspace*{4.2cm}

{\large\bfseries\color{lhblue}En bref}\par\smallskip
\begin{tabularx}{\linewidth}{@{}>{\bfseries}l X@{}}
\toprule
Objectif & Faire entrer des photos dans le Lakehouse et les analyser, automatiser tout le pipeline avec
           Airflow, l'optimiser, l'exposer aux équipes métier avec Trino et Superset, puis le surveiller
           (Prometheus, Grafana) et l'observer (OpenLineage, Marquez). \\
Durée    & 3 séances de 3 h (+ environ 3 h de travail personnel). \\
Prérequis & TP 1 et TP 2 terminés (tables Silver, projet dbt). Images du cours en version
            \cmd{2026.2} (section~\ref{sec:deroule}). \\
À rendre & Un compte rendu, les notebooks exécutés, le dossier \cmd{pipeline/}, les deux DAG et le
           tableau de bord Superset (section~\ref{sec:livrables}). \\
\bottomrule
\end{tabularx}

% =====================================================================
\section{Le scénario}
% =====================================================================

Olist ouvre un \textbf{catalogue photo} : chaque vendeur envoie les photos de ses produits, avec une
fiche descriptive. La direction pose trois exigences :

\begin{enumerate}[leftmargin=*,nosep]
  \item \textbf{Fiabilité :} aucune photo illisible, minuscule ou en double ne doit arriver dans le
        catalogue, et les fiches mal catégorisées doivent être repérées.
  \item \textbf{Automatisation :} le traitement doit tourner seul, chaque jour, et s'arrêter de
        lui-même si un lot est douteux.
  \item \textbf{Accès :} les équipes métier veulent interroger les données en SQL et disposer d'un
        tableau de bord partagé, sans dépendre de l'équipe data.
\end{enumerate}

Ce TP termine le module : vous assemblez une architecture complète, de la donnée brute au tableau de
bord, comme en entreprise.

% =====================================================================
\section{Les données}
% =====================================================================

\begin{center}
\small
\renewcommand{\arraystretch}{1.25}
\begin{tabularx}{\textwidth}{@{}l l X r@{}}
\toprule
\textbf{Source} & \textbf{Format} & \textbf{Contenu} & \textbf{Volume} \\
\midrule
Photos des produits & JPEG & Une photo par produit, 256 pixels au plus, rangées par vendeur
  (\cmd{data/catalogue/photos/vendeur=XX/}) & $\approx$ 3\,000 fichiers \\
Fiches produits & JSON & Une fiche par ligne : titre, type de produit, marque, couleur, vendeur
  (\cmd{data/catalogue/produits.jsonl}) & 2\,906 fiches \\
Modèle de vision & ONNX & Le réseau MobileNet, pré-entraîné (\cmd{data/modeles/}) & 3,7 Mo \\
Tables des TP 1 et 2 & Iceberg & Commandes, avis, livraisons, schéma en constellation & -- \\
\bottomrule
\end{tabularx}
\end{center}

Les photos et les fiches sont un sous-ensemble du jeu public \textbf{Amazon Berkeley Objects}
(Amazon et université de Berkeley, licence CC BY-NC 4.0) : six types de produits (chaises,
canapés, tables, tapis, chaussures et sacs à main), environ 500 par type, répartis entre douze
vendeurs fictifs. Comme dans la
réalité, le lot contient des défauts, ajoutés volontairement : fichiers coupés pendant l'envoi,
fichier qui n'est pas une image, vignettes minuscules, même photo envoyée deux fois, ou recadrée.

\begin{notion}
\textbf{Une image est une donnée non structurée :} pour un ordinateur, ce n'est qu'une grille de
pixels. Pour l'exploiter, on en extrait des informations \emph{structurées} : ses dimensions, une
empreinte qui permet de reconnaître les doublons, et un \textbf{vecteur} de caractéristiques
calculé par un réseau de neurones, qui résume ce qu'elle montre.
\end{notion}

% =====================================================================
\section{L'architecture de bout en bout}
% =====================================================================

La figure~\ref{fig:archi} rassemble tout ce que vous avez construit pendant le module. Le TP 3 y
ajoute les photos, l'\textbf{orchestration} (Airflow, qui déclenche et surveille chaque étape) et
l'\textbf{exposition} (Trino et Superset, pour les équipes métier), puis la \textbf{supervision} et
l'\textbf{observabilité} de toute la plateforme.

\begin{figure}[H]
\centering
\resizebox{\textwidth}{!}{%
\begin{tikzpicture}[
  font=\sffamily\small, >=Stealth,
  col/.style={draw=#1,thick,rounded corners=4pt,fill=#1!7,text width=3.3cm,align=left,
              inner sep=6pt,font=\sffamily\footnotesize,minimum height=3.0cm},
  op/.style={font=\sffamily\scriptsize,text=black!70,align=center},
  num/.style={circle,fill=lhblue,text=white,font=\sffamily\bfseries\scriptsize,inner sep=1.4pt}]

\node[col=gray] (src) at (0,0) {\textbf{Sources}\\[3pt] PostgreSQL (TP 1)\\ CSV, avis (TP 1, 2)\\ JSON du suivi (TP 2)\\ \textbf{photos, fiches (TP 3)}};
\node[col=darkgray] (land) at (4.3,0) {\textbf{Landing}\\ \texttt{s3://landing/}\\[3pt] fichiers déposés\\ tels quels};
\node[col=lhbronze] (br) at (8.6,0) {\textbf{Bronze}\\[3pt] copie fidèle\\ \texttt{photos} : octets\\ et empreinte \texttt{sha256}};
\node[col=lhsilver] (si) at (12.9,0) {\textbf{Silver}\\[3pt] contrôlé, dédoublonné\\ \texttt{photos}\\ + quarantaine};
\node[col=lhgold] (go) at (17.2,0) {\textbf{Gold}\\[3pt] constellation (dbt)\\ \texttt{photos\_vecteurs}\\ \texttt{photos\_controle}};

\draw[->,thick] (src) -- node[num]{1} (land);
\draw[->,thick] (land) -- node[num]{2} (br);
\draw[->,thick] (br) -- node[num]{3} (si);
\draw[->,thick] (si) -- node[num]{4} (go);

% Orchestration
\node[draw=lhblue,thick,rounded corners=4pt,fill=lhblue!5,minimum width=17.6cm,minimum height=1.1cm,
      font=\sffamily\footnotesize,align=center] (af) at (10.75,3.0)
  {\textbf{Apache Airflow} -- DAG \texttt{catalogue\_photos} (chaque jour) et \texttt{maintenance\_iceberg} (chaque semaine)};
\node[num] at ([xshift=4mm]af.west) {5};
\node[draw=lhteal,thick,rounded corners=4pt,fill=lhteal!6,minimum width=17.6cm,minimum height=0.8cm,
      font=\sffamily\footnotesize,align=center] (obs) at (10.75,4.45)
  {\textbf{Supervision et observabilité} -- Prometheus + Grafana (métriques, alertes) \quad OpenLineage + Marquez (lineage)};
\foreach \n in {land,br,si,go} \draw[->,lhblue,thick] (af.south -| \n) -- (\n.north);

% Exposition et usages
\node[col=lhblue,minimum height=1.4cm,text width=3.3cm] (tr) at (17.2,-3.6) {\textbf{Trino}\\ SQL interactif\\ et fédération};
\node[col=lhblue,minimum height=1.4cm,text width=3.3cm] (ss) at (12.9,-3.6) {\textbf{Superset}\\ tableau de bord\\ de l'entreprise};
\node[col=lhblue,minimum height=1.4cm,text width=3.3cm] (ml) at (21.5,-3.6) {\textbf{Vision}\\ MobileNet (ONNX)\\ + scikit-learn};
\draw[->,thick] (go.south) -- node[num]{6} (tr.north);
\draw[->,thick] (tr.west) -- node[num]{7} (ss.east);
\draw[<->,thick] (go.east) -| node[num,pos=0.75]{8} (ml.north);
\draw[->,thick] (src.south) -- ++(0,-4.1) -| node[num,pos=0.25]{9} (tr.south);
\node[op,anchor=south west] at (0.2,-5.0) {fédération : Trino lit aussi PostgreSQL en direct};
\end{tikzpicture}}
\caption{L'architecture complète du module. Les numéros renvoient aux étapes décrites ci-dessous.}
\label{fig:archi}
\end{figure}

\begin{description}[leftmargin=1.4em,labelindent=0pt,itemsep=3pt,font=\normalfont]
\item[\etape{1}] \textbf{Dépôt :} les vendeurs envoient leurs photos dans la zone d'atterrissage,
  par lots (script \cmd{deposer\_photos.py}).
\item[\etape{2}] \textbf{Bronze :} chaque nouvelle photo devient une ligne de \cmd{bronze.photos} :
  son chemin, son vendeur, sa taille, son empreinte exacte (\cmd{sha256}) et ses octets.
\item[\etape{3}] \textbf{Silver :} chaque photo est décodée et contrôlée ; les photos illisibles,
  trop petites ou en double partent en quarantaine.
\item[\etape{4}] \textbf{Gold :} le réseau MobileNet calcule le vecteur de chaque photo ; un
  classifieur reconnaît le type de produit et signale les fiches incohérentes. dbt reconstruit
  la constellation du TP 2.
\item[\etape{5}] \textbf{Orchestration :} Airflow enchaîne ces étapes chaque jour, attend les
  nouvelles photos, relance une étape en cas d'erreur passagère, arrête tout si un contrôle
  échoue, et compacte les tables chaque semaine.
\item[\etape{6}] \textbf{SQL interactif :} Trino lit les tables Gold, sans Spark et sans copie.
\item[\etape{7}] \textbf{Tableau de bord :} Superset présente les indicateurs, à partir de Trino.
\item[\etape{8}] \textbf{Vision par ordinateur :} les vecteurs servent à classer les photos et à
  retrouver les photos semblables.
\item[\etape{9}] \textbf{Fédération :} Trino joint, dans une même requête, la base PostgreSQL de
  l'application de vente et les tables du Lakehouse.
\end{description}

% =====================================================================
\section{Déroulé}\label{sec:deroule}
% =====================================================================

\textbf{Mettre à jour l'environnement (une fois, au début de la séance 1) :} le TP 3 utilise de
nouvelles versions des images (avec la lecture d'images et Airflow). Dans le terminal Ubuntu, depuis
le dossier du dépôt :
\begin{lstlisting}[style=shell]
git pull
cp .env.example .env          # nouvelles versions des images (2026.2)
docker compose --profile full --profile orchestration --profile observabilite pull
docker compose up -d
\end{lstlisting}
Si vous aviez modifié \cmd{.env} (machine de 8 Go), refaites la même modification. Puis, dans un
terminal JupyterLab, copiez les énoncés :
\begin{lstlisting}[style=shell]
cp -r notebooks/tp3 notebooks/mon_travail/
\end{lstlisting}

\begin{center}
\small
\renewcommand{\arraystretch}{1.25}
\begin{tabularx}{\textwidth}{@{}l X l@{}}
\toprule
\textbf{Partie} & \textbf{Contenu} & \textbf{Services nécessaires} \\
\midrule
1 & Les photos dans le Lakehouse (notebook \nb{01\_images\_catalogue}) & de base \\
2 & Le pipeline et Airflow (\nb{02\_pipeline\_airflow}, \nb{pipeline/}, \nb{dags/}) & \cmd{-{}-profile orchestration} \\
3 & Mesurer et optimiser (\nb{03\_optimiser}) & de base, Airflow \\
4 & Trino et Superset (\nb{04\_trino\_superset}) & \cmd{-{}-profile full} \\
5 & Surveiller et observer (\nb{05\_surveiller\_observer}) & \cmd{-{}-profile observabilite}, Airflow \\
\bottomrule
\end{tabularx}
\end{center}

\begin{attention}
\textbf{Mémoire :} JupyterLab, Airflow, Trino et Superset ne doivent pas tourner tous ensemble sur une
machine de 8 Go. Arrêtez ce qui ne sert pas à la partie en cours :
\cmd{docker compose -{}-profile orchestration stop airflow} ou
\cmd{docker compose -{}-profile full stop trino superset}. Les outils de la partie 5 ajoutent
environ 1,5 Go.
\end{attention}

% =====================================================================
\section{Partie par partie}
% =====================================================================

\subsection*{Partie 1 -- Des photos dans le Lakehouse \hfill \nb{01\_images\_catalogue}}

Vous ingérez le premier lot de photos en Bronze, vous contrôlez chaque photo en Silver, puis vous
calculez les vecteurs MobileNet en Gold, pour retrouver les photos semblables et contrôler les
types de produits déclarés.

\begin{notion}
\textbf{Deux sortes de doublons :} l'empreinte \cmd{sha256} identifie un fichier \emph{à l'octet
près} : elle reconnaît la même photo, même renommée ou envoyée par un autre vendeur. Mais une photo
recadrée ou recompressée a une empreinte complètement différente : on applique donc aussi une règle
métier, \og une seule photo par produit \fg{}. Comparer les photos d'un produit à l'autre serait
trompeur : beaucoup de fiches sont des variantes d'un même modèle, photographiées de la même façon.
\end{notion}

\begin{notion}
\textbf{L'apprentissage par transfert :} MobileNet a appris, sur des millions d'images, à reconnaître
des formes, des textures et des parties d'objets. On ne le réentraîne pas : on récupère, pour chaque
photo, le vecteur de 576 nombres qu'il calcule avant sa décision finale, et une simple régression
logistique apprend nos six types de produits à partir de ces vecteurs. Le modèle est fourni au
format \textbf{ONNX}, lisible sans PyTorch.
\end{notion}

\begin{resultat}
Le lot 1 compte 1\,977 photos ; 76 partent en quarantaine (22 trop petites, 21 doublons exacts,
17 illisibles, 16 photos en double). Le type de produit est reconnu sur environ 95~\% des photos ;
une cinquantaine de fiches sont signalées, dont beaucoup sont réellement mal catégorisées (un canapé
déclaré comme chaise, des barquettes déclarées comme table).
\end{resultat}

\capture{tp3_01_doublons}{Partie 1 : un produit envoyé avec deux photos, la même recadrée.}{fig:res01a}
\capture{tp3_01_similaires}{Partie 1 : une photo et les cinq photos les plus semblables selon MobileNet.}{fig:res01b}

\subsection*{Partie 2 -- Du notebook au pipeline \hfill \nb{02\_pipeline\_airflow}}

Le code du notebook devient quatre scripts, un par étape, puis un DAG Airflow les enchaîne
(figure~\ref{fig:dag}).

\begin{figure}[H]
\centering
\begin{tikzpicture}[font=\sffamily\footnotesize, >=Stealth, node distance=0.4cm,
  t/.style={draw=lhblue,thick,rounded corners=3pt,fill=lhblue!6,align=center,inner sep=4pt,minimum height=0.95cm,text width=1.95cm}]
\node[t] (a) {\texttt{attendre\_}\\\texttt{photos}};
\node[t,right=of a] (b) {\texttt{bronze\_}\\\texttt{photos}};
\node[t,right=of b] (c) {\texttt{silver\_}\\\texttt{photos}};
\node[t,right=of c] (d) {\texttt{gold\_}\\\texttt{photos}};
\node[t,right=of d] (e) {\texttt{controle\_}\\\texttt{qualite}};
\node[t,right=of e] (f) {\texttt{dbt\_build}};
\foreach \x/\y in {a/b,b/c,c/d,d/e,e/f} \draw[->,thick] (\x) -- (\y);
\node[font=\sffamily\scriptsize,text=black!65,below=2pt of a,align=center] {capteur};
\node[font=\sffamily\scriptsize,text=lhred,below=2pt of e,align=center] {bloquant};
\end{tikzpicture}
\caption{Le DAG \texttt{catalogue\_photos}. Une tâche ne démarre que si la précédente a réussi.}
\label{fig:dag}
\end{figure}

\begin{notion}
\textbf{Orchestrer un pipeline :} un orchestrateur comme \textbf{Apache Airflow} décrit les étapes
d'un traitement et leur ordre dans un \textbf{DAG}, un graphe de tâches écrit en Python. Il les
déclenche selon un \textbf{calendrier} (\cmd{@daily}) ou quand une condition est remplie (un
\emph{capteur} attend de nouvelles photos), \textbf{relance} une tâche après une erreur passagère,
n'exécute une tâche que si les précédentes ont réussi, et garde le \textbf{journal} et la durée de
chaque exécution.
\end{notion}

\begin{notion}
\textbf{Idempotence :} une étape idempotente peut être relancée sans changer le résultat. C'est
indispensable dès qu'un pipeline est automatisé : les relances sont fréquentes (panne, relance
automatique, rejeu). Les scripts du TP sont incrémentaux : ils ne traitent que les photos qu'ils
n'ont pas encore vues.
\end{notion}

\begin{resultat}
Le DAG s'exécute sans erreur sur le lot 1, puis, après le dépôt du lot 2, ne traite que les
nouvelles photos. Avec un seuil de rejet abaissé à 1~\%, la tâche \cmd{controle\_qualite} échoue
après ses relances, et \cmd{dbt\_build} n'est pas exécutée.
\end{resultat}

\capture{tp3_02_graph}{Partie 2 : le DAG \texttt{catalogue\_photos} dans Airflow, vue \emph{Graph}.}{fig:res02a}
\capture{tp3_02_gantt}{Partie 2 : la durée de chaque tâche, vue \emph{Gantt}, pour le lot 1 et le lot 2.}{fig:res02b}

\subsection*{Partie 3 -- Mesurer, puis optimiser \hfill \nb{03\_optimiser}}

Vous mesurez l'effet des petits fichiers, vous compactez une table, vous partitionnez une table par
vendeur, puis vous confiez la maintenance à un second DAG, hebdomadaire.

\begin{notion}
\textbf{Ranger les fichiers :} chaque écriture ajoute des fichiers et une version à une table
Iceberg. Le \textbf{compactage} (\cmd{rewrite\_data\_files}) regroupe les petits fichiers ; la
\textbf{purge} (\cmd{expire\_snapshots}) supprime les versions trop anciennes. Le
\textbf{partitionnement} range les fichiers par valeur d'une colonne : une requête filtrée ne lit
que les fichiers utiles.
\end{notion}

\begin{resultat}
Après 90 petites écritures, la table de démonstration compte plus de 90 fichiers ; le compactage
les ramène à un seul. La table partitionnée compte un fichier par vendeur.
\end{resultat}

\capture{tp3_03_fichiers}{Partie 3 : le nombre de fichiers avant et après le compactage.}{fig:res03}

\subsection*{Partie 4 -- Trino et Superset \hfill \nb{04\_trino\_superset}}

Vous interrogez le Lakehouse avec Trino, vous joignez la base PostgreSQL et les tables Iceberg dans
une même requête, puis vous construisez le tableau de bord de l'entreprise dans Superset.

\begin{notion}
\textbf{Séparer le stockage et le calcul :} les tables Iceberg sont de simples fichiers dans S3,
décrits par le catalogue. Plusieurs moteurs les utilisent \emph{en même temps} et sans copie :
Spark pour les traitements lourds, Trino pour les requêtes interactives, Superset au-dessus de
Trino. Chacun se dimensionne et se remplace indépendamment des données.
\end{notion}

\begin{resultat}
La requête fédérée montre les quelques commandes de PostgreSQL absentes de \cmd{silver.orders} :
celles mises en quarantaine au TP 1.
\end{resultat}

\subsection*{Partie 5 -- Surveiller et observer \hfill \nb{05\_surveiller\_observer}}

Tout est déjà configuré : vous utilisez les outils comme une équipe d'exploitation. Vous lisez le
tableau de bord Grafana pendant une exécution du DAG, vous créez une alerte sur les tâches en échec,
puis vous suivez dans Marquez le lineage de \cmd{gold.photos\_controle}.

\begin{attention}
Démarrez la supervision \emph{avant} Airflow, ou redémarrez Airflow après elle
(\cmd{docker compose -{}-profile orchestration restart airflow}) : Airflow ne cherche le service de
métriques qu'à son démarrage. Le message \cmd{Could not configure StatsClient} dans ses logs signifie
qu'il n'enverra aucune métrique ; il est normal pendant les parties 2 à 4, où la supervision n'est pas
lancée.
\end{attention}

\begin{center}
\small
\renewcommand{\arraystretch}{1.25}
\begin{tabularx}{\textwidth}{@{}l X X@{}}
\toprule
& \textbf{Monitoring (supervision)} & \textbf{Observabilité des données} \\
\midrule
Question & Les \emph{services} tournent-ils bien ? & Les \emph{données} sont-elles fiables, d'où viennent-elles ? \\
Mesures & processeur, mémoire, durée et échecs des tâches & lineage, historique des exécutions, schémas \\
Outils du TP & Prometheus (métriques) + Grafana (tableaux de bord, alertes) & OpenLineage (format des événements) + Marquez (stockage, interface) \\
En entreprise & Datadog, CloudWatch, Azure Monitor & DataHub, Unity Catalog, Monte Carlo, Atlan \\
\bottomrule
\end{tabularx}
\end{center}

\begin{notion}
\textbf{Le lineage :} le graphe qui relie chaque table aux jobs qui la produisent et aux tables
qu'ils lisent. Avec OpenLineage, chaque job Spark envoie lui-même la description de ses entrées et de
ses sorties à chaque exécution : le graphe n'est pas dessiné à la main, il est \emph{déduit des
exécutions réelles}. Il sert à trouver l'origine d'une erreur (\emph{analyse de cause}) et à savoir
ce qui sera touché par un changement (\emph{analyse d'impact}).
\end{notion}

\begin{center}
\small
\renewcommand{\arraystretch}{1.2}
\begin{tabularx}{\textwidth}{@{}l l X@{}}
\toprule
\textbf{Outil} & \textbf{Adresse} & \textbf{Connexion} \\
\midrule
Grafana    & \url{http://localhost:3000} & \cmd{admin} / \cmd{admin}, tableau \emph{Lakehouse -- supervision} \\
Prometheus & \url{http://localhost:9090} & sans mot de passe (onglet \emph{Graph} pour essayer PromQL) \\
Marquez    & \url{http://localhost:3001} & sans mot de passe, espace de noms \cmd{lakehouse} \\
\bottomrule
\end{tabularx}
\end{center}

\capture{tp3_05_grafana}{Partie 5 : le tableau de bord Grafana pendant une exécution du DAG.}{fig:grafana}
\capture{tp3_05_lineage}{Partie 5 : le lineage de \texttt{gold.photos\_controle} dans Marquez.}{fig:marquez}

% =====================================================================
\section{Le tableau de bord}\label{sec:dashboard}
% =====================================================================

Dans Superset (\url{http://localhost:8088}, \cmd{admin} / \cmd{admin}) :
\begin{enumerate}[leftmargin=*,itemsep=2pt]
  \item \textbf{Connecter Trino :} \emph{Settings > Database Connections > + Database}, choisir
        \emph{Trino}, adresse \cmd{trino://superset@trino:8080/lakehouse}, puis \emph{Test connection}.
  \item \textbf{Déclarer les jeux de données :} \emph{Datasets > + Dataset} pour
        \cmd{gold.fait\_avis}, \cmd{gold.fait\_livraisons} et \cmd{gold.photos\_controle}. Dans
        \emph{SQL Lab}, exécutez la requête de l'étape 4 du notebook 04 et enregistrez-la comme jeu
        de données (\emph{Save dataset}).
  \item \textbf{Construire au moins cinq graphiques :} la note moyenne et la part d'avis négatifs
        (indicateurs), la part d'avis négatifs par mois, les incidents par transporteur, la part de
        photos rejetées par vendeur, et le tableau des fiches incohérentes.
  \item \textbf{Assembler le tableau de bord}, avec un \textbf{filtre} commun (par exemple l'état
        du client), et donnez-lui un titre clair.
\end{enumerate}

\capture{tp3_dashboard}{Le tableau de bord Superset attendu.}{fig:dash}

% =====================================================================
\section{Livrables}\label{sec:livrables}
% =====================================================================

\subsection*{Le compte rendu}

Il ne recopie pas le code : il explique vos choix et vos résultats.
\begin{enumerate}[leftmargin=*,itemsep=2pt]
  \item \textbf{Architecture} : votre schéma de bout en bout, et le rôle de chaque composant.
  \item \textbf{Images} : le stockage choisi, les règles de qualité, le bilan de la quarantaine, et ce
        que montrent la recherche de photos semblables et le contrôle des types.
  \item \textbf{Orchestration} : votre DAG, l'idempotence de vos scripts, le comportement en cas
        d'échec et de relance, les captures des vues \emph{Graph} et \emph{Gantt}.
  \item \textbf{Optimisation} : vos mesures avant et après, et vos choix de partitionnement.
  \item \textbf{Exposition} : la requête fédérée, une capture du tableau de bord Superset, et la
        comparaison avec Streamlit.
  \item \textbf{Supervision et observabilité} : ce que montrent Grafana et Marquez, l'alerte créée,
        et la différence entre les deux approches.
  \item \textbf{Réponses aux questions} des notebooks, rédigées et argumentées.
  \item \textbf{Limites et améliorations} : ce qu'il faudrait de plus pour passer en production
        (droits d'accès, montée en charge, coûts, astreinte).
\end{enumerate}

\clearpage
% =====================================================================
\section*{Annexe -- Aide-mémoire}\label{sec:aidememoire}
% =====================================================================

\begin{center}
\small
\renewcommand{\arraystretch}{1.22}
\begin{tabularx}{\textwidth}{@{}>{\raggedright\arraybackslash}p{4.1cm} >{\raggedright\arraybackslash\ttfamily\footnotesize}X@{}}
\toprule
\normalfont\textbf{Je veux\dots} & \normalfont\textbf{Code} \\
\midrule
\multicolumn{2}{@{}l}{\textbf{\textcolor{lhblue}{Fichiers et images}}} \\
Lire des fichiers binaires & spark.read.format("binaryFile").option("recursiveFileLookup", True).load(chemin) \\
Empreinte exacte d'un fichier & F.sha2("content", 256) \\
Décoder une image & Image.open(io.BytesIO(octets)) \quad img.load() \\
\midrule
\multicolumn{2}{@{}l}{\textbf{\textcolor{lhblue}{Airflow}}} \\
Démarrer / arrêter Airflow & docker compose -{}-profile orchestration up -d \quad \dots stop airflow \\
Une tâche qui lance un script & BashOperator(task\_id="x", bash\_command="python x.py") \\
Attendre une condition & PythonSensor(task\_id="x", python\_callable=f, mode="reschedule") \\
Ordre des tâches & a >{}> b >{}> c \\
Relances automatiques & default\_args=\{"retries": 2, "retry\_delay": timedelta(minutes=1)\} \\
Calendrier & schedule="@daily" \quad schedule="0 3 * * 0" \textrm{(dimanche, 3 h)} \\
\midrule
\multicolumn{2}{@{}l}{\textbf{\textcolor{lhblue}{Iceberg}}} \\
Fichiers d'une table & SELECT * FROM lakehouse.silver.photos.files \\
Compacter & CALL lakehouse.system.rewrite\_data\_files(table => 'silver.photos') \\
Purger l'historique & CALL lakehouse.system.expire\_snapshots(table => '...', older\_than => TIMESTAMP '...') \\
Partitionner & CREATE TABLE ... USING iceberg PARTITIONED BY (vendeur) AS SELECT ... \\
\midrule
\multicolumn{2}{@{}l}{\textbf{\textcolor{lhblue}{Trino}}} \\
Nom complet d'une table & lakehouse.gold.fait\_avis \quad postgres.ventes.orders \\
Versions d'une table & SELECT * FROM lakehouse.silver."orders\$snapshots" \\
Relire une version & SELECT ... FROM lakehouse.silver.orders FOR VERSION AS OF <snapshot\_id> \\
Connexion Superset & trino://superset@trino:8080/lakehouse \\
\midrule
\multicolumn{2}{@{}l}{\textbf{\textcolor{lhblue}{Supervision}}} \\
Démarrer les outils & docker compose -{}-profile observabilite -{}-profile orchestration up -d \\
Puis redémarrer Airflow & docker compose -{}-profile orchestration restart airflow \\
Mémoire des conteneurs (PromQL) & container\_memory\_working\_set\_bytes\{name=\~{}"lakehouse-.*"\} \\
Tâches en échec sur 5 min & sum(increase(airflow\_taches\_en\_echec\_total[5m])) \\
\bottomrule
\end{tabularx}
\end{center}

\end{document}
